\subsection*{Partie I}
\begin{enumerate}
 \item D'après les définitions :
\begin{displaymath}
 \Mat_{\mathcal B}=
\begin{bmatrix}
 0      & 0      & \cdots & \cdots & 0      & -a_0     \\
 1      & 0      &        &        &        & -a_1     \\
 0      & 1      & \ddots &        & \vdots & \vdots   \\
 \vdots & 0      & \ddots & \ddots &        & \vdots   \\
 \vdots & \vdots &        & \ddots & 0      & -a_{p-2} \\
 0      & 0      & \cdots & 0      & 1      & -a_{p-1} \\
\end{bmatrix}
\end{displaymath}
\item On a bien $Q(f)\circ f = f\circ Q(f)$ car $f$ commute avec ses puissances.\newline
D'après la définition, pour tout $i<p$, $f^i(e_1)=e_{i+1}$ et
\begin{displaymath}
 f^p(e_1)=f(e_p)= -a_0 e_1 -a_1e_2 - \cdots - a_{p-1}e_p
\end{displaymath}
On en déduit
\begin{displaymath}
 P(f)(e_1)= a_0 e_1 + a_1e_2 + \cdots + a_{p-1}e_p + f(e_p) = 0_E 
\end{displaymath}
On vérifie alors que $P(f)(e_i)$ est nul pour tous les $i$ entre $2$ et $P$ car
\begin{displaymath}
 P(f)(e_i)=P(f)\circ f^{i-1}(e_1)=f^{i-1}\circ P(f)(e_1)=f^{i-1}(O_E)=0_E
\end{displaymath}
L'endomorphisme $P(f)$ est donc nul puisqu'il prend la valeur $0_E$ sur tous les vecteurs d'une base.
\item \begin{enumerate}
 \item Soit $\lambda$ une valeur propre de $f$. Il existe alors un vecteur $u\neq 0_E$ tel que $f(u)=\lambda u$. On en déduit que $f^i(u)=\lambda^i u$ pour tous les entiers naturels $i$. Cela entraîne
\begin{displaymath}
 0_E=P(f)(u)=P(\lambda)u
\end{displaymath}
d'où $P(\lambda)=0$ car $u$ est non nul.
\item D'après les conditions de l'énoncé, $Q$ est un polynôme unitaire de degré $n-1$. Il existe donc des complexes $b_0,\cdots, b_{p-2}$ tels que
\begin{align*}
 Q =& b_O + b_1 X + \cdots +b_{p-2}X^{p-2}+X^{p-1}\\
Q(f)(e_1) =& b_Oe_1 + b_1 e_2 + \cdots +b_{p-2}e_{p-1}+e_{p}
\end{align*}
Comme la famille $(e_1,\cdots,e_p)$ est libre et que 
\begin{displaymath}
 (b_0,\cdots, b_{p-2},1) \neq (0,0,\cdots,0)
\end{displaymath}
le vecteur $Q(f)(e_1)$ est non nul. Notons $v=Q(f)(e_1)$ ce vecteur non nul. Alors:
\begin{multline*}
 P=(X-\lambda)Q \Rightarrow (f-\lambda \Id)\circ Q(f) =O_{\mathcal L(E)}
\Rightarrow (f-\lambda \Id)\circ Q(f)(e_1) =O_E \\
\Rightarrow f(v)=\lambda v
\end{multline*}
Le complexe $\lambda$ est donc bien une valeur propre.
\end{enumerate}
\end{enumerate}

\subsection*{Partie II}
\begin{enumerate}
\item Considérons un $\C$-espace vectoriel $F$ quelconque de dimension $p$ muni d'une base $\mathcal U =(u_1,\cdots,u_p)$.\\ Définissons un endomorphisme $\varphi$ de $F$ par sa matrice dans la base $\mathcal U$ :
\begin{displaymath}
 \Mat_{\mathcal U} f = A
\end{displaymath}
D'après le cours, $A$ est inversible si et seulement si $f$ est bijective. Comme $f$ est un endomorphisme d'un espace vectoriel de dimension finie, $f$ est bijective si et seulement si elle est injective. On en déduit que $A$ est non inversible si et seulement si $f$ est non injective c'est à dire de noyau non nul. \`A chaque vecteur non nul du noyau correspond la colonne $X$ de ses coordonnées dans la base $\mathcal U$. Cette colonne vérifie $AX=0$.
\item Soit $A$ une matrice non inversible et
\begin{displaymath}
 X=\begin{bmatrix}
    x_1\\
x_2\\
\vdots \\
x_p
   \end{bmatrix}
\end{displaymath}
une colonne telle que $AX=0$. C'est à dire
\begin{align*}
 \forall i\in\{1,\cdots,p\} & & \sum_{k=1}^{p}a_{ik}x_k = 0
\end{align*}
Il existe un entier $m$ entre 1 et $p$ tel que 
\begin{displaymath}
 |x_m|=\max (|x_1|,\cdots,|x_p|)
\end{displaymath}
\'Ecrivons la relation pour $i=m$ et formons une inégalité
\begin{multline*}
 -a_{mm}x_m = \sum_{k\in\{1,\cdots,p\}-\{m\}}a_{mk}x_k
\Rightarrow
|a_{mm}||x_m|\leq \sum_{k\in\{1,\cdots,p\}-\{m\}}|a_{mk}||x_k|\\
\Rightarrow
|a_{mm}||x_m|\leq \sum_{k\in\{1,\cdots,p\}-\{m\}}|a_{mk}||x_m|=r_m(A)|x_m| \\
\Rightarrow |x_m|\left( |a_{mm}| - r_m(A)\right) \leq 0 
\Rightarrow |a_{mm}|\leq r_m(A)
\end{multline*}
car $X$ étant non nulle, $|x_m|>0$.
\item Si $\lambda$ est une valeur propre de $A$, la matrice $A-\lambda I_p$ n'est pas inversible. Il existe donc un $m$ tel que 
\begin{displaymath}
 |a_{mm}-\lambda|\leq r_m(A-\lambda I_p)=r_m(A)
\end{displaymath}
Ceci signifie que $\lambda$ est dans le $m$ ieme disque de  Ger\v{s}gorin. Toutes les valeurs propres sont donc dans le domaine de  Ger\v{s}gorin de $A$.
\item Soit $P\in\C[X]$, on a vu en I.3. que les racines de $P$ sont les valeurs propres de $A$. Elles sont donc dans le domaine de  Ger\v{s}gorin de $A$.\newline
Précisons ce domaine dans le cas où $P=-1+iX-4X^2-(1+i)X^3+X^4$.
\begin{figure}
 \centering
 \includegraphics[width=5cm]{Cgersg_1.pdf}
 % Cgersg_1.pdf: 1179666x1179666 pixel, 0dpi, infxinf cm, bb=
 \caption{Domaine de Ger\v{s}gorin pour $P=-1+iX-4X^2-(1+i)X^3+X^4$}
 \label{fig: Cgersg_1}
\end{figure}

\begin{displaymath}
 A=
\begin{bmatrix}
 0 & 0 & 0 &  1 \\
 1 & 0 & 0 &  -i \\
 0 & 1 & 0 &  4 \\
 0 & 0 & 1 &  1+i \\
\end{bmatrix}
\Rightarrow
\left\lbrace 
\begin{aligned}
 r_1(A) =& 1 \\
 r_2(A) =& 2 \\
 r_3(A) =& 5 \\
 r_4(A) =& 1 \\
\end{aligned}
\right. 
\end{displaymath}
Les trois premiers disques sont centrés à l'origine. Ils sont tous dans le disque $\mathcal D$ centré à l'origine et de rayon 5. Le dernier disque est centré au point d'affixe $1+i$ et de rayon 1. Il est aussi contenu dans le disque $\mathcal D$.\\
Précisons ce domaine dans le cas où $P=-2+iX-jX^2-4X^3+X^4$.
\begin{figure}
 \centering
 \includegraphics[width=6cm]{Cgersg_2.pdf}
 % Cgersg_1.pdf: 1179666x1179666 pixel, 0dpi, infxinf cm, bb=
 \caption{Domaine de Ger\v{s}gorin pour $P=-2+iX-jX^2-4X^3+X^4$}
 \label{fig: Cgersg_2}
\end{figure}

\begin{displaymath}
 A=
\begin{bmatrix}
 0 & 0 & 0 &  2 \\
 1 & 0 & 0 &  -i \\
 0 & 1 & 0 &  j \\
 0 & 0 & 1 &  4 \\
\end{bmatrix}
\Rightarrow
\left\lbrace 
\begin{aligned}
 r_1(A) =& 2 \\
 r_2(A) =& 2 \\
 r_3(A) =& 2 \\
 r_4(A) =& 1 \\
\end{aligned}
\right. 
\end{displaymath}
Les trois premiers disques sont confondus centrés à l'origine et de rayon 2. Le dernier disque est centré au point d'affixe $4$ et de rayon 1.
\end{enumerate}
